\documentclass[11pt]{article}

\usepackage[T1]{fontenc}
\usepackage[utf8]{inputenc}
\usepackage{lmodern}
\usepackage{amsmath,amssymb,amsthm,mathtools}
\usepackage{booktabs}
\usepackage{enumitem}
\usepackage[margin=1in]{geometry}
\usepackage{microtype}
\usepackage{xcolor}
\usepackage[hidelinks]{hyperref}

\newtheorem{theorem}{Theorem}[section]
\newtheorem{lemma}[theorem]{Lemma}
\newtheorem{proposition}[theorem]{Proposition}
\newtheorem{corollary}[theorem]{Corollary}
\theoremstyle{definition}
\newtheorem{definition}[theorem]{Definition}
\newtheorem{algorithm}[theorem]{Algorithm}
\theoremstyle{remark}
\newtheorem{remark}[theorem]{Remark}

\newcommand{\Sn}{S_n}
\newcommand{\TD}{\operatorname{TD}}
\newcommand{\dist}{d_t}
\newcommand{\E}{\mathcal E}
\newcommand{\D}{\mathcal D}
\newcommand{\T}{\mathcal T}
\newcommand{\CanSort}{\mathsf{CanSort}}
\newcommand{\canon}{\operatorname{can}}
\newcommand{\Stab}{\operatorname{Stab}}
\newcommand{\id}{\iota}
\newcommand{\code}[1]{\texttt{#1}}

% This switch is intentionally true until the exhaustive degree-18 run and
% its independent verification have been completed and archived.
\newif\ifdraftstatus
\draftstatustrue

\title{An Exact Deletion--Extension Approach to the\\
  Block-Transposition Diameter of $S_{18}$}
\author{Author list and affiliations to be finalized}
\date{Working paper, \today}

\begin{document}
\maketitle

\begin{abstract}
The block-transposition diameter $\TD(n)$ is the largest minimum number of
adjacent-block swaps required to sort a permutation in $S_n$.  The known value
$\TD(17)=10$ and the one-symbol deletion inequality imply
$10\leq\TD(18)\leq11$.  We develop an exact deletion--extension computation
which decides between these two values without enumerating the distance-ten
layer in $S_{18}$.  Complete exact layers in degree $n-1$ are expanded through
their full extended-toric orbits, the new maximum symbol is inserted in every
position, and every surviving canonical child is resolved by deletion bounds
or exhaustive bounded search.  The search uses the odd-cycle lower bound,
distance-preserving contractions, a slack-exact move filter, and a generalized
root prepass based on all standardized one-symbol deletions.  Two further root
reductions use qualifying odd-cycle covers and a newly verified $N=19$
support-clearing classification.  We state the
mathematical completeness conditions, describe the persistent and GPU
implementation, and report exhaustive small-degree audits and controlled
microbenchmarks.  Completion of the distance-eleven search with an empty
result proves $\TD(18)=10$; a surviving child instead proves
$\TD(18)=11$.
\end{abstract}

\ifdraftstatus
\begin{center}
\fcolorbox{red!60!black}{red!4}{%
  \begin{minipage}{0.91\linewidth}
  \textbf{Computational status.}
  This document is a proof-ready working paper, not yet a claim that
  $\TD(18)=10$ has been established.  The reduction, implementation, audits,
  and benchmark are present; the complete degree-18 run and its archived
  audit artifacts remain pending.  No degree-18 count, running time, hash, or
  empty result is inferred from a partial computation.
  \end{minipage}}
\end{center}
\fi

\section{Introduction}

A block transposition exchanges two adjacent, nonempty blocks in a
permutation.  Sorting by block transpositions was introduced in its modern
form by Bafna and Pevzner~\cite{BafnaPevzner1998}; its diameter problem asks
for the greatest distance from the identity in the corresponding Cayley
graph.  Eriksson et al.~\cite{ErikssonEtAl2001} established the general upper
bound
\[
  \TD(n)\leq \left\lfloor\frac{2n-2}{3}\right\rfloor
  \qquad (n\geq9)
\]
and introduced toric equivalence as a central symmetry of the problem.
Korchmaros supplied a detailed algebraic treatment of the toric maps and
their distance invariance~\cite{Korchmaros2015,Korchmaros2014}.  Korchmaros
and Kov\'acs later identified the dihedral factor in the automorphism group
of the block-transposition Cayley graph~\cite{KorchmarosKovacs2017}.
Those results justify the root quotient and the orbit--stabilizer weights used
here.  They do not, without an additional state-canonicalization argument,
justify quotienting the internal nodes of a bounded search.

The exact value $\TD(16)=9$ was recently obtained by Silva et
al.~\cite{SilvaEtAl2026}, giving the presently known sequence through degree
17 and, in particular, $\TD(17)=10$.  Degree 18 is therefore the next
undecided case.  Deleting any one symbol and standardizing changes distance
by at most one.  It follows immediately that
\[
  10=\TD(17)\leq\TD(18)\leq\TD(17)+1=11.       \tag{1.1}
\]
Consequently, the problem is an exact emptiness decision: does any
permutation in $S_{18}$ have distance 11?

B\'ona and Pittel~\cite{BonaPittel2025} recently counted broad families with
restricted cycle lengths and used them to obtain many permutations with a
large block-transposition lower bound.  Their enumerative result does not by
itself exclude distance 11 in $S_{18}$, but its cycle bound is exactly the
admissible lower bound used by our DFS.  The underlying 3-cycle length formula
and cover characterization are due to M\"uhle and
Nadeau~\cite{MuhleNadeau2019}.  We extract a further algorithmic consequence
from those facts in Proposition~\ref{prop:slack}: when the remaining budget
equals the lower bound, only a move increasing the relevant odd-cycle count by
two can occur on a successful path.

The central idea of this paper is to generate that possible upper layer from
the complete degree-17 diametral layer, while quotienting by all
extended-toric symmetries.  This is substantially smaller than a scan of
$18!$, but symmetry reduction must be used carefully.  Insertion does not
commute with choosing a single canonical parent.  We therefore expand every
stored parent through all distinct ordinary members of its symmetry orbit
\emph{before} inserting the new maximum.  Canonicalization is applied only to
the resulting children.

The second contribution is an exact bounded-search engine tailored to these
children.  The odd-cycle lower bound determines which move deltas can still
fit in the remaining budget.  A complete degree-$(n-1)$ layer also supports a
root deletion prepass for \emph{every} degree $n$: one sufficiently easy
deletion proves that the root is sortable within budget, whereas one deletion
whose exact distance exceeds the budget disproves it.  Neither optimization
is a degree-18 special case.

Our intended final theorem has the following computational premise.

\begin{theorem}[Target theorem]
\label{thm:target}
If the exhaustive computation specified in Sections~\ref{sec:algorithm}
and~\ref{sec:certificate} completes with no member of $\E(18,11)$, and the
completion conditions in Theorem~\ref{thm:certificate} are verified, then
$\TD(18)=10$.
\end{theorem}

The theorem is presently conditional only because the final exhaustive run
has not been completed.  All mathematical reductions apply independently of
the hardware and implementation.

\section{Definitions and conventions}
\label{sec:definitions}

We write permutations in one-line notation
$p=[p_1\ p_2\ \cdots\ p_n]\in S_n$ and write
$\id=[1\ 2\ \cdots\ n]$.  For cuts
\[
  0\leq a<b<c\leq n,
\]
the block transposition $\tau_{a,b,c}$ sends $p$ to
\[
 [p_1\cdots p_a\mid
   p_{b+1}\cdots p_c\mid
   p_{a+1}\cdots p_b\mid
   p_{c+1}\cdots p_n].
\]
There are $\binom{n+1}{3}$ such moves.  The block-transposition distance
$\dist(p)$ is the minimum number of moves required to transform $p$ into
$\id$, and
\[
  \TD(n)=\max_{p\in S_n}\dist(p).
\]

We use three related layer notations:
\begin{align*}
  \E(n,d)&=\{p\in S_n:\dist(p)=d\},\\
  \D(n)&=\E(n,\TD(n)),\\
  \T(n,k)&=\{p\in S_n:\dist(p)\geq k\}.
\end{align*}
The computation persists exact layers $\E(n,d)$, not merely lower-bound
candidate sets.

For a symbol $x$ occurring in $p$, let $p-x$ denote the word obtained by
deleting $x$ and then standardizing the remaining values to
$\{1,\ldots,n-1\}$.  Every reference to deletion below includes this
standardization.

For the cycle formulation let $N=n+1$, let
\[
  \sigma=(0\ 1\ \cdots\ n),\qquad
  \beta_p=(0\ p_1\ \cdots\ p_n),
\]
and use left composition: $(fg)(x)=f(g(x))$.  Two products appear in the
implementation:
\[
  \omega_p=\sigma\beta_p^{-1},
  \qquad
  \gamma_p=\beta_p^{-1}\sigma.
\]
They are conjugate, since $\omega_p=\sigma\gamma_p\sigma^{-1}$, and hence
have the same cycle type.  The canonical key is the image vector of
$\omega_p$; the lower-bound kernel traverses $\gamma_p$.  Keeping these two
coordinate conventions distinct prevents an otherwise easy cut-indexing
error.

\section{Deletion, insertion, and exact layers}
\label{sec:deletion}

The following inequality is the basis of both the layer extension and the
new root prepass.  It is Lemma~2.4 of Ara\'ujo et
al.~\cite{AraujoEtAl2026}.

\begin{lemma}[One-symbol deletion inequality]
\label{lem:deletion}
For every $p\in S_n$ and every symbol $x$ of $p$,
\[
  \dist(p-x)\leq\dist(p)\leq\dist(p-x)+1.       \tag{3.1}
\]
\end{lemma}

The first inequality is restriction monotonicity: deleting a symbol from a
sorting sequence leaves a valid sequence after empty moves are suppressed.
The second is the one-symbol insertion bound.  In the usual adjacency proof,
one preliminary block transposition places the inserted symbol in a
distance-preserving adjacency, after which a sorting sequence of the deletion
lifts to the original permutation.  We use Lemma~\ref{lem:deletion} as a
published combinatorial result rather than as a property of the program.

\begin{corollary}[Deletion interval]
\label{cor:interval}
For $p\in S_n$,
\[
  \max_x\dist(p-x)\leq\dist(p)
  \leq 1+\min_x\dist(p-x).                      \tag{3.2}\label{eq:deletion-interval}
\]
In particular, the distances of two one-symbol deletions of $p$ differ by at
most one.
\end{corollary}

If the deletion distances include both $q$ and $q+1$, the interval in
\eqref{eq:deletion-interval} collapses and $\dist(p)=q+1$.  If every deletion has distance
$q$, only $q$ and $q+1$ remain possible, so one exact bounded decision is
necessary.

\begin{proposition}[Exact-layer recurrence]
\label{prop:layer-recurrence}
Every member of $\E(n,d)$ is obtained by inserting the maximum symbol $n$
into a member of
\[
  \E(n-1,d-1)\cup\E(n-1,d).                    \tag{3.3}\label{eq:layer-recurrence}
\]
Conversely, an insertion from the right-hand side belongs to $\E(n,d)$
exactly when its deletion interval and, when necessary, an exhaustive
$\CanSort$ decision assign it distance $d$.
\end{proposition}

\begin{proof}
Delete the maximum symbol from $p\in\E(n,d)$.  By
Lemma~\ref{lem:deletion}, the deletion has distance $d-1$ or $d$, proving
coverage.  For the converse, compute all deletion distances.  If the interval
in Corollary~\ref{cor:interval} is a singleton, it is the exact distance.  In
the remaining case all deletions have the same distance $q$, and exhaustive
search decides whether the child is sortable in $q$ moves.  Failure gives
distance $q+1$ by the upper bound in Lemma~\ref{lem:deletion}.
\end{proof}

An equivalent threshold formulation is often convenient.

\begin{proposition}[Threshold recurrence]
\label{prop:threshold}
The set $\T(n,k)$ can be constructed from maximum insertions into
$\T(n-1,k-1)$ by applying the following tests to each child $p$:
\begin{enumerate}[label=(\roman*)]
  \item reject $p$ if any deletion is outside $\T(n-1,k-1)$;
  \item accept $p$ if any deletion belongs to $\T(n-1,k)$;
  \item otherwise every deletion has distance $k-1$, and accept $p$ exactly
    when exhaustive search proves $\neg\CanSort(p,k-1)$.
\end{enumerate}
\end{proposition}

\begin{proof}
If a deletion has distance at most $k-2$, Lemma~\ref{lem:deletion} gives
$\dist(p)\leq k-1$.  If a deletion has distance at least $k$, restriction
monotonicity gives $\dist(p)\geq k$.  The only unresolved case is the uniform
distance-$(k-1)$ case, where the insertion bound leaves precisely the two
possibilities $k-1$ and $k$.
\end{proof}

This recurrence explains why the complete diametral layer alone is
insufficient when a diameter remains constant.  For example, generating
$\D(n)$ may also require the degree-$(n-1)$ layer immediately below
$\D(n-1)$.  The implementation therefore requests exact layers lazily by
\eqref{eq:layer-recurrence}; it does not assume that every step increases the diameter.
This distinction is important at degree 14: the published ordinary count
$|\D(14)|=5{,}246{,}800{,}005$~\cite[Table~5.2]{Korchmaros2015} already
precludes treating a full unquotiented output list as a convenient
intermediate representation.

\subsection{The degree-18 reduction}

Monotonicity and Lemma~\ref{lem:deletion} give
\[
  \TD(n-1)\leq\TD(n)\leq\TD(n-1)+1.
\]
Since $\TD(17)=10$, the only possible degree-18 upper layer is
$\E(18,11)$.  If $p$ belongs to that layer, every one-symbol deletion has
distance 10: it has distance at least 10 by Lemma~\ref{lem:deletion}, and at
most $\TD(17)=10$.  In particular, deletion of 18 lies in $\D(17)$.

\begin{corollary}[Diameter-only decision]
\label{cor:diameter-only}
The complete set of maximum insertions into $\D(17)$ contains every member of
$\E(18,11)$.  Therefore
\[
  \E(18,11)=\varnothing\quad\Longleftrightarrow\quad\TD(18)=10,
\]
whereas any member of $\E(18,11)$ proves $\TD(18)=11$.
\end{corollary}

The computation need not construct $\E(18,10)$ to decide the diameter.

\section{Extended-toric symmetry}
\label{sec:symmetry}

Toric maps preserve block-transposition distance
\cite{ErikssonEtAl2001,Korchmaros2014,Korchmaros2015}.  Adjoining the reverse
map gives a dihedral action of order $2N$ on the fixed-target instances; this
is the factor identified in the Cayley-graph automorphism theorem of
Korchmaros and Kov\'acs~\cite{KorchmarosKovacs2017}.

Let $\mathfrak r(x)=-x\pmod N$, and define $g^h=hgh^{-1}$.  With the target
$\sigma$ fixed, the implementation compares the image vectors
\[
  \omega^{\sigma^v}
  \quad\hbox{and}\quad
  (\omega^{-1})^{\sigma^{v+1}\mathfrak r},
  \qquad 0\leq v<N.                             \tag{4.1}\label{eq:dihedral-action}
\]
The canonical representative $\canon(p)$ is the lexicographically least
image vector among these $2N$ transforms.  If $G_n$ denotes this dihedral
action, a canonical key represents
\[
  w(p)=\frac{2N}{|\Stab_{G_n}(\omega_p)|}        \tag{4.2}\label{eq:orbit-weight}
\]
ordinary permutations.  The stabilizer correction is essential; simply
assigning weight $2N$ would give incorrect counts for symmetric instances.

\begin{proposition}[Distance preservation]
\label{prop:distance-preservation}
All ordinary permutations represented by one canonical key have the same
block-transposition distance.
\end{proposition}

\begin{proof}
The rotations in \eqref{eq:dihedral-action} are the fixed-target form of the toric maps,
and the second family composes these maps with reversal.  Both induce
automorphisms of the block-transposition Cayley graph fixing the identity up
to the fixed-target coordinate change.  Hence they preserve distance.
\end{proof}

Canonicalizing a parent and immediately inserting into only that stored word
is not complete: insertion positions are not invariant under the action.
The following order of operations is therefore mandatory.

\begin{proposition}[Orbit expansion before insertion]
\label{prop:orbit-expansion}
Suppose that one canonical representative of every required parent orbit in
$S_{n-1}$ is stored.  Expand each representative into all distinct ordinary
members of its $G_{n-1}$-orbit, insert $n$ in every one of the $n$ positions,
and retain only canonical children.  Then no $G_n$-orbit of eligible children
is lost.
\end{proposition}

\begin{proof}
Let $q$ be an eligible child, and let $p=q-n$.  The store contains a canonical
representative $p_0$ in the orbit of $p$.  Full orbit expansion of $p_0$
therefore emits $p$.  Inserting $n$ at its original position emits $q$, after
which canonicalization retains the representative of the orbit of $q$.
\end{proof}

Moreover, each ordinary child has a unique ordinary parent obtained by
deleting its maximum.  Expanding distinct parent orbits into distinct
ordinary parents therefore avoids a global ordinary-child deduplication
table.  Canonical keys and orbit weights supply the compact persisted form.

\section{The exact extension algorithm}
\label{sec:algorithm}

For a requested exact layer $\E(n,d)$, the dependency planner first ensures
that $\E(n-1,d-1)$ and $\E(n-1,d)$ are complete whenever those layers are
within the established diameter bounds.  It then executes
Algorithm~\ref{alg:extension}.

\begin{algorithm}[Canonical exact-layer extension]
\label{alg:extension}
Given complete canonical stores for the required parent layers:
\begin{enumerate}[label=\arabic*.]
  \item Stream every canonical parent and recover all its distinct ordinary
    extended-toric orbit members.
  \item Insert the new maximum $n$ in all $n$ positions of every member.
  \item Discard a child unless its fixed-target $\omega$ key is canonical.
  \item Standardize every one-symbol deletion and query its canonical key in
    the complete parent stores or resident exact tables.
  \item Use Corollary~\ref{cor:interval} to determine the child's distance
    whenever possible.  If all deletion distances are equal to $q$, run the
    exhaustive decision $\CanSort(p,q)$.
  \item Persist the canonical child exactly when its distance is the requested
    $d$; add its orbit weight only to aggregate counts.
\end{enumerate}
\end{algorithm}

For the diameter-only degree-18 run, the requested layer is $\E(18,11)$.
The program stops early on a verified distance-eleven child.  If none exists,
it must process the complete parent manifest and settle every canonical child
before it can publish an empty result.

\subsection{Candidate intervals and exact decisions}

Let
\[
  m(p)=\min_x\dist(p-x),\qquad
  M(p)=\max_x\dist(p-x).
\]
Corollary~\ref{cor:interval} gives
\[
  M(p)\leq\dist(p)\leq m(p)+1.                  \tag{5.1}
\]
Since $M(p)-m(p)\leq1$, a mixed pair of deletion distances determines
$\dist(p)=M(p)$.  Only $m(p)=M(p)=q$ reaches bounded search.  The search is a
decision procedure, not a heuristic: it enumerates every block transposition
unless a proved lower bound, exact table lookup, or distance-preserving
contraction settles the state.

\subsection{Lazy dependency schedule}

Exact stores are horizon-independent.  A request for $\E(n,d)$ recursively
requests only $\E(n-1,d-1)$ and $\E(n-1,d)$.  This matters in the constant
diameter steps at degrees 14 and 16, where the layer below the previous
diametral set is required.  For the degree-18 diameter-only decision,
$\E(17,11)$ is outside the known upper bound and $\E(17,10)=\D(17)$ is the
only parent layer at the final step.

\section{Exact bounded search}
\label{sec:search}

The bounded search receives a root $p\in S_n$ and a budget $r$.  It performs
iterative depth-first search over all $\binom{n+1}{3}$ cut triples.  The only
unconditional move suppression is immediate reversal of the previous move,
which cannot occur in a shortest path.  GPU kernel quanta may suspend and
resume a search, but a quantum boundary never classifies an unfinished root.

\subsection{Odd-cycle lower bound}

Let $c_{\mathrm{odd}}(p)$ be the number of odd cycles of
$\gamma_p=\beta_p^{-1}\sigma$.  The Bafna--Pevzner lower bound is
\[
  h(p)=\frac{N-c_{\mathrm{odd}}(p)}{2}
  \leq\dist(p).                                 \tag{6.1}
\]
This is precisely the 3-cycle length formula
$\ell_3(\gamma_p)=(N-c_{\mathrm{odd}}(p))/2$ of M\"uhle and
Nadeau~\cite[Proposition~3.1]{MuhleNadeau2019}; its role in large-distance
enumeration is restated by B\'ona and
Pittel~\cite[Section~4]{BonaPittel2025}.

A block transposition changes $c_{\mathrm{odd}}$ by a value
\[
  \Delta\in\{-2,0,2\}.
\]
In the applicable-3-cycle representation, M\"uhle and
Nadeau~\cite[Proposition~3.2 and Corollary~3.3]{MuhleNadeau2019}
characterize exactly when right multiplication changes $\ell_3$ by
$-1$, $0$, or $1$, equivalently when it traverses a cover.  Thus an
applicable move has exactly the three displayed odd-cycle deltas, and a
$+2$ move traverses a cover toward the identity.
The production kernel computes this delta using the same constant-time cycle
routine as the independently audited emitter.  Its cut arguments are formed
in the $\gamma_p$ convention, rather than copied from the canonical
$\omega_p$ convention.

\begin{proposition}[Slack-exact move filter]
\label{prop:slack}
At a search state with remaining budget $r$, let $s=r-h(p)$.  A child can fit
within budget only under the following rule:
\[
\begin{array}{c|c}
  s & \text{admissible }\Delta\\ \hline
  0 & 2\\
  1 & 0,2\\
  s\geq2 & -2,0,2.
\end{array}                                      \tag{6.2}\label{eq:slack-filter}
\]
Thus filtering the other moves preserves completeness.
\end{proposition}

\begin{proof}
Since the child has
$h'=h-\Delta/2$ and remaining budget $r-1$, its slack is
\[
  s'=(r-1)-h'=s-1+\frac{\Delta}{2}.              \tag{6.3}
\]
A necessary condition is $s'\geq0$.  Substitution of the three possible
values of $\Delta$ gives exactly \eqref{eq:slack-filter}.
\end{proof}

This rule is a direct consequence of the M\"uhle--Nadeau length formula and
cover characterization, not an assumption about the structure of optimal
sorting sequences.

\subsection{Qualifying covers and contraction}

A $+2$ move is called \emph{qualifying} when it creates a fixed point in the
resulting algebraic permutation.  The fixed point gives an identity adjacency
in the word coordinates, so the child contracts by one symbol without changing
distance.  This is the qualifying-move mechanism used in the proof of
$\TD(16)=9$~\cite[Lemma~7.2]{SilvaEtAl2026}.

\begin{proposition}[Root qualifying-move prepass]
\label{prop:qualifying-prepass}
Let $p\in S_n$ have remaining budget $r$.  If an applicable $+2$ move produces
a child whose new fixed point contracts to $q\in S_{n-1}$, and a complete
exact lookup proves $\dist(q)\leq r-1$, then $\CanSort(p,r)$.
Failure to find such a move or lookup is inconclusive.
\end{proposition}

\begin{proof}
Contraction preserves the child's distance, and the qualifying move costs one,
so $\dist(p)\leq1+\dist(q)\leq r$.
\end{proof}

The resident-table implementation scans all qualifying moves.  The
persisted-store implementation emits the contracted canonical key for only the
first qualifying move and queries the complete degree-$(n-1)$ store; a failed
or inconclusive first lookup leaves the root to exact DFS.  This asymmetry can
lose an optimization opportunity but cannot change a decision.

\subsection{The \texorpdfstring{$N=19$}{N=19} oriented-unicycle reduction}

Silva et al.~\cite[Proposition~7.4]{SilvaEtAl2026} show for the odd-cycle
lengths needed in their $N=17$ proof that an oriented odd cycle either admits
the qualifying criterion or has a support-clearing sequence of length
$(k+1)/2$, apart from one explicitly handled local exception.  We extended
their exhaustive route-O classification to $N=19$.

\begin{lemma}[$N=19$ unicycle classification]
\label{lem:n19-clearing}
Let $(\omega,\beta)$ be a realizable cyclic-target pair on 19 symbols and
suppose that $\omega$ is an oriented 19-cycle.  If it evades the qualifying
criterion of Silva et al., then it has a support-clearing sequence of ten block
transpositions.
\end{lemma}

\begin{proof}[Computational verification]
Normalize $\beta=(0\ 1\ \cdots\ 18)$ and write the cycle uniquely from
zero.  The exhaustive route-O enumeration visited $9{,}304{,}845$ partial
nodes and retained three normalized cycles.  Removing the unoriented survivor
and quotienting by the auxiliary dihedral action leaves two oriented orbit
representatives:
\begin{verbatim}
(0 10 1 11 2 12 3 13 4 14 5 15 6 16 7 17 8 18 9)
(0 16 13 10 7 4 1 17 14 11 8 5 2 18 15 12 9 6 3)
\end{verbatim}
Both representatives pass the cyclic-target realizability test, and the native
witness search verifies a ten-move support-clearing sequence for each.  The
enumeration prunes only conditions necessary for failure of the qualifying
criterion, so these representatives cover every unresolved oriented case.
\end{proof}

At a degree-18 DFS root, fixed points are contracted before this reduction is
tested.  Consequently an unresolved odd unicycle is fixed-point-free and its
nontrivial cycle has length 19, so Lemma~\ref{lem:n19-clearing} supplies the
ten-move upper bound.  The reduction is enabled only for $n=18$ and budgets at
least 10; otherwise the exact DFS continues.

\subsection{A generalized deletion prepass}

The deletion inequality is also useful \emph{inside} every bounded search.
The following statement is independent of the target degree.

\begin{proposition}[Root deletion prepass]
\label{prop:root-prepass}
Let $p\in S_n$ and let $r$ be a bounded-search budget.  For any standardized
one-symbol deletion $q=p-x$:
\begin{enumerate}[label=(\roman*)]
  \item an upper bound $\dist(q)\leq r-1$ proves
    $\CanSort(p,r)$;
  \item a lower bound $\dist(q)>r$ proves
    $\neg\CanSort(p,r)$.
\end{enumerate}
If neither condition holds for any deletion, the result is inconclusive and
the complete DFS must continue.
\end{proposition}

\begin{proof}
Part (i) follows from the insertion half of
Lemma~\ref{lem:deletion}:
$\dist(p)\leq\dist(q)+1\leq r$.  Part (ii) follows from restriction
monotonicity: $\dist(p)\geq\dist(q)>r$.
\end{proof}

In the implementation, the prepass is enabled whenever a degree-$(n-1)$
exact pattern table or completed membership bitmap is resident.  An exact
entry supplies both bounds.  For a completed top-layer bitmap for
$\E(n-1,t)$, a set bit means exact distance $t$, while an unset bit means
distance at most $t-1$; a partial bitmap is never queried.  All $n$
standardized deletions are checked at DFS depth zero.  The code is
parameterized by $n$ and contains no branch specific to $n=18$.

The prepass may settle a large fraction of roots when a complete
$\D(n-1)$ layer is resident, but it does not change the dependency plan: it
reuses a layer already loaded for extension or contraction.  Its absence or
an inconclusive answer affects performance only.

\subsection{Distance-preserving contraction and exact tables}

If a word contains an identity adjacency---including the two boundary
adjacencies---the later symbol can be deleted without changing distance.
Repeated contraction produces a core permutation.  A resident exact table
can then settle the core.  A table entry records an exact distance; an absent
entry below a declared retained floor records only the corresponding upper
bound.  These semantics are kept separate in the device encoding.

\section{Implementation and persistence}
\label{sec:implementation}

The reference implementation is in the \code{experiments} module of the
project repository~\cite{SourceCode}.  Java coordinates persistent storage,
host queues, and CUDA devices.  CUDA kernels perform orbit expansion,
maximum insertion, child canonicality, deletion ranking, resident table
lookups, and exact DFS.  GPU mode never silently falls back to a CPU sorting
decision.

\subsection{Canonical records}

For a degree-$n$ instance, the key is the canonical image vector of
$\omega=\sigma\beta^{-1}$.  Every symbol occupies five bits, most significant
bit first.  With $n=18$, the 19 images occupy 12 bytes.  The value is one byte
whose low nibble is the exact distance.  Orbit weights are recomputed during
orbit expansion and are stored only as aggregate manifest counts.

The persistent backend is RocksDB~\cite{RocksDB}.  Each exact layer has its
own directory and manifest.  A readable manifest records the degree, exact
distance, encoding version, canonical count, weighted ordinary count, and
the property \code{complete=true}.  A store without a compatible complete
manifest is unavailable, never an empty mathematical set.

\subsection{Atomic completion}

Records are written asynchronously, but the completion manifest is published
only after workers have terminated successfully, the writer queue has drained,
the database has flushed, and expected counts have been checked when known.
The diameter result is first written to a temporary file and then atomically
renamed.  An exception, interruption, incomplete parent traversal, or
unresolved candidate therefore cannot publish a completed empty result.

For an empty diameter decision, the program verifies three coverage
equalities:
\begin{align*}
  \text{parents processed}
    &=\text{canonical size in the input manifest},\\
  \text{ordinary children generated}
    &=n\sum_d\text{weighted parent count at distance }d,\\
  \text{canonical candidates}
    &=\text{deletion-rejected}+
      \text{directly determined}+
      \text{exact-search decisions}.            \tag{7.1}\label{eq:coverage}
\end{align*}
The second equality incorporates full orbit expansion and all maximum
insertion positions.

\subsection{GPU caches}

Completed canonical stores may be expanded into one of two device caches.
A four-bit dense table can hold exact distances for several retained layers;
a one-bit membership table represents one complete exact layer.  Population
is checked against the source manifest after orbit expansion.  All ranks,
offsets, and populations use 64-bit counters.  Device-specific cache
selection respects reserved search workspaces and a configurable VRAM
fraction.  Failure to fit a cache causes a slower lookup path, not a change in
the mathematical search.

When neither exact representation fits, the implementation may instead load
the union of the required complete canonical parent layers into a GPU Bloom
filter.  The filter is used only in the negative direction: a definite miss
rejects a child during deletion checking, whereas every possible match is
confirmed by an exact lookup in RocksDB.  False positives therefore add I/O
but cannot alter a distance decision.  The loader checks the number of source
records against the manifests and re-queries every loaded key before enabling
the filter, thereby auditing the absence of false negatives in the constructed
device image.

\subsection{Parallel work and early termination}

Canonical candidates are aggregated by search budget and distributed over
independent GPU slots.  Full queues apply backpressure.  A positive
distance-eleven result cancels the remaining decision work, because one
witness settles the diameter.  Negative publication, in contrast, requires
complete traversal and the coverage equalities in \eqref{eq:coverage}.
Sparse unresolved roots are retained until either a useful device batch is
available or the extension input is exhausted.  A final underfilled batch is
partitioned among the available streams, allowing its independent serial DFS
roots to occupy multiple multiprocessors.  Kernel quanta bound the time between
cancellation checks; their size, final partitioning, and the minimum batch
threshold are performance parameters and do not affect which DFS states are
visited.

\section{Correctness of the computational reduction}
\label{sec:correctness}

\begin{theorem}[Exactness of a completed extension]
\label{thm:extension-correct}
Assume that:
\begin{enumerate}[label=(\alph*)]
  \item every input exact-layer manifest is complete and correct;
  \item each stored parent is expanded into its full distinct ordinary
    extended-toric orbit;
  \item all maximum insertion positions are generated;
  \item deletion queries use complete exact stores with the stated absence
    semantics;
  \item every unresolved candidate is decided by a complete bounded search;
  \item every root prepass uses only the proved one-sided conclusions in
    Propositions~\ref{prop:root-prepass} and
    \ref{prop:qualifying-prepass}, or the complete classification in
    Lemma~\ref{lem:n19-clearing};
  \item output is published only after all coverage equalities hold.
\end{enumerate}
Then the published output for a request $(n,d)$ contains exactly one canonical
representative of every extended-toric orbit in $\E(n,d)$, with correct orbit
weights.
\end{theorem}

\begin{proof}
For completeness, Proposition~\ref{prop:layer-recurrence} places the
maximum deletion of every $p\in\E(n,d)$ in one of the requested parent
layers.  Proposition~\ref{prop:orbit-expansion} shows that full orbit
expansion followed by all insertions generates the orbit of $p$ and retains
its canonical key.  Corollary~\ref{cor:interval} and exhaustive bounded
search assign its exact distance, so it is retained.

For soundness, every retained key comes from a valid insertion.  Mixed
deletion distances determine the exact distance by
Corollary~\ref{cor:interval}; a uniform deletion distance leaves only two
possibilities, which complete bounded search distinguishes.  Distance is
constant on each canonical orbit by
Proposition~\ref{prop:distance-preservation}.  Orbit--stabilizer gives the
weight in \eqref{eq:orbit-weight}.  Finally, the maximum deletion is unique for each
ordinary child, and the coverage checks rule out a truncated traversal.
\end{proof}

\begin{theorem}[Degree-18 completion certificate]
\label{thm:certificate}
Suppose a run requesting $\E(18,11)$ satisfies all hypotheses of
Theorem~\ref{thm:extension-correct}, reads the archived complete
$\E(17,10)$ input, finishes without early cancellation, verifies
\eqref{eq:coverage}, and atomically publishes an empty result.  Then
$\TD(18)=10$.
\end{theorem}

\begin{proof}
Theorem~\ref{thm:extension-correct} makes the published empty set equal to
$\E(18,11)$.  Corollary~\ref{cor:diameter-only} then gives
$\TD(18)=10$.
\end{proof}

If the run instead stops on a child for which exhaustive search proves
$\neg\CanSort(p,10)$, its maximum deletion has distance 10 and hence
$\dist(p)\leq11$.  Therefore $\dist(p)=11$, proving the alternative value
$\TD(18)=11$.

\section{Validation}
\label{sec:validation}

Small-degree validation tests implementation behavior; it does not replace
the degree-18 exhaustive run.  The current validation suite contains the
following independent or cross-implementation checks.

\begin{enumerate}
  \item Independent breadth-first search reconstructs complete distance
    histograms and exact sets through degree 9 (and degree 10 in the extended
    audit), and verifies the deletion inequality on every permutation.
  \item CPU and GPU fixed-target canonicality and orbit weights are compared
    on every permutation through degree 9.  The weighted representatives sum
    to $n!$.
  \item Complete GPU diametral counts through degree 9 are checked against the
    independent histograms, including runs whose largest exact table has
    degree $n-1$.
  \item The actual emitter cycle-delta functions are compared with explicit
    cycle recounting for every permutation and every oriented cut triple
    through degree 8, together with a nontrivial degree-9 merge case.
  \item The production lower-bound coordinate mapping and
    Proposition~\ref{prop:slack} are audited for every permutation/cut pair
    through degree 8.  This is 3,696,986 queries in total.
  \item Forced one-, two-, and three-tick GPU kernel quanta test resumability,
    root exhaustion, and both phases of the two-budget protocol.
  \item Queue overflow is forced and checked to split the input range without
    losing or double-counting a canonical root.
  \item Completed device membership and pattern tables are population-checked
    against persistent manifests after full orbit expansion.
  \item A persisted-$\D(9)$ integration run reconstructs
    $|\D(10)|=255{,}053$ and resolves 494 roots through the qualifying-move
    lookup without changing the exact output.
  \item The $N=19$ route-O enumeration visits $9{,}304{,}845$ nodes and
    leaves the two realizable oriented orbits in
    Lemma~\ref{lem:n19-clearing}; native CUDA witness search and the
    production prepass independently resolve both in ten moves.
\end{enumerate}

The production slack audit independently constructs each child word, rebuilds
its $\gamma$ permutation, recounts cycles, and compares this result with both
the constant-time delta routine and the feasibility rule in \eqref{eq:slack-filter}.
This specifically protects the distinction between the $\omega$ and
$\gamma$ coordinate systems.

\section{Benchmark of the slack and deletion rules}
\label{sec:benchmark}

We measured four variants of the same CUDA bounded-search kernel: neither
rule, only Proposition~\ref{prop:slack}, only
Proposition~\ref{prop:root-prepass}, and both rules.  Each run used 32,768
evenly spaced permutation ranks, a degree-$(n-1)$ resident exact table, two
warmups, and the median of eleven measured repetitions on an NVIDIA RTX 4070
Laptop GPU.  Every variant returned an identical decision array.

\begin{table}[ht]
\centering
\caption{Kernel time and speedup relative to the unfiltered bounded DFS.}
\label{tab:benchmark}
\begin{tabular}{@{}rrrrrrrrr@{}}
\toprule
$n$ & budget & baseline & \multicolumn{2}{c}{slack only}
 & \multicolumn{2}{c}{deletions only} & \multicolumn{2}{c}{both}\\
\cmidrule(lr){4-5}\cmidrule(lr){6-7}\cmidrule(l){8-9}
 & & ms & ms & speedup & ms & speedup & ms & speedup\\
\midrule
8  & 4 &  6.089 &  2.854 & 2.133 &  5.765 & 1.056 &  2.666 & 2.283\\
9  & 4 &  0.975 &  0.715 & 1.364 &  0.626 & 1.558 &  0.288 & 3.388\\
10 & 5 & 30.952 & 13.117 & 2.360 & 30.401 & 1.018 & 13.030 & 2.375\\
\bottomrule
\end{tabular}
\end{table}

These are bounded-search microbenchmarks, not projections of the degree-18
wall time.  Root distributions, cache residency, I/O, orbit expansion, and
the fraction settled before DFS all change with degree.  The benchmark shows
that both rules are exact and reusable, and that their interaction can be
substantial; it does not estimate the time needed to prove
$\E(18,11)=\varnothing$.

The benchmark is reproducible with
\begin{verbatim}
bash experiments/run.sh search-benchmark --benchmark-n 9 \
  --benchmark-budget 4 --benchmark-jobs 32768 \
  --benchmark-table-degree 8 --benchmark-repeats 11
\end{verbatim}
These two features default to enabled and can be disabled independently with
the Java properties \code{extension.slackMoveFilter} and
\code{extension.rootDeletionPrepass}.  The qualifying prepass is controlled by
\code{extension.rootMovePrepass} (default \code{auto}, with minimum budget 8),
and the degree-18 clearing reduction by
\code{extension.oddCyclePrepass}.  They are covered by the exact integration
tests above rather than folded into this earlier four-way timing table.

\section{Final computation and proof artifacts}
\label{sec:certificate}

The intended diameter-only invocation builds native CUDA kernels, opens the
complete dependency stores, and requests the upper layer at degree 18.  A
representative command is
\begin{verbatim}
mvn -pl experiments -am -Pcuda-native -DgpuTests=true verify \
  dependency:copy-dependencies
bash experiments/run.sh diameter --no-build --max-n 18 \
  --cuda-ptx-dir "$PWD/experiments/target/native" \
  --data-dir experiments/data/extension \
  --log-file experiments/logs/extension.log
\end{verbatim}
Machine-specific GPU lists, slot counts, and VRAM caps may be added without
changing the searched set.

For publication, the final run should archive more than the small atomic
result properties file.  At minimum, the archive should contain:
\begin{enumerate}
  \item the exact source revision and a clean-worktree patch status;
  \item compiler, JDK, Maven, CUDA toolkit, driver, JCuda, and RocksDB versions;
  \item build commands and SHA-256 hashes of all loaded PTX files;
  \item complete input manifests and hashes for every exact store;
  \item the full command line, environment settings, hardware inventory, and
    start/end timestamps;
  \item an uninterrupted log containing the parent, generated, canonical,
    rejected, determined, and exact-query totals;
  \item the atomic \code{td-18.properties} result and its hash;
  \item output from all CPU and GPU audits at the archived revision;
  \item preferably, an independent verifier which checks the result without
    invoking the production search scheduler.
\end{enumerate}

\begin{table}[ht]
\centering
\caption{Fields to replace from the completed degree-18 run.  Pending entries
must not be filled from a partial run.}
\label{tab:final-run}
\begin{tabular}{@{}p{0.34\linewidth}p{0.57\linewidth}@{}}
\toprule
Field & Archived value\\
\midrule
Source revision & \code{PENDING}\\
Input $\E(17,10)$ manifest/hash & \code{PENDING}\\
Hardware and software stack & \code{PENDING}\\
Canonical parents processed & \code{PENDING}\\
Ordinary insertions generated & \code{PENDING}\\
Canonical candidates settled & \code{PENDING}\\
Bounded-search decisions & \code{PENDING}\\
Start, end, and wall time & \code{PENDING}\\
Outcome for $\E(18,11)$ & \code{PENDING}\\
Result/log archive hashes & \code{PENDING}\\
Independent verification & \code{PENDING}\\
\bottomrule
\end{tabular}
\end{table}

\subsection{Interpreting the two outcomes}

If a match is found, the archive must include its ordinary word and canonical
$\omega$ key.  An independent bounded search can then verify that it is not
sortable in ten moves, while a degree-17 deletion supplies the upper bound
11.  If no match is found, the log and manifests must show that every input
parent and canonical child was settled and that no early-stop flag was set.
Only the latter complete negative outcome activates
Theorem~\ref{thm:certificate}.

\section{Limitations and further certificate work}

The persistent layer stores contain canonical permutations and exact
distances, but not sorting sequences.  A completed negative search is thus a
reproducible computational proof based on source, manifests, and audit logs,
not a compact independently replayable certificate.  For submission, a
separate checker would materially strengthen the result.  It should have a
minimal interface to the emitter, independently implement the combinatorial
case split, and verify either all emitted bounded-search certificates or a
more compact exhaustive-search transcript.  The emitter/checker separation in
Silva et al.~\cite{SilvaEtAl2026} provides a useful model.

The current validation degrees are deliberately small enough for independent
BFS.  Agreement there is strong regression evidence, but it does not prove
that a degree-18 run completed.  Likewise, benchmark speedups establish only
performance and decision agreement on the sampled roots.

Finally, the extended-toric quotient is used only at the root/class level.
The dihedral Cayley-graph automorphisms do not, by themselves, justify
discarding noncanonical states inside one representative's DFS.  The present
DFS therefore explores all cut triples subject only to the proved rules in
Section~\ref{sec:search}.

\section{Conclusion}

The deletion inequality reduces the open degree-18 diameter to the emptiness
of a single exact layer, $\E(18,11)$.  Extended-toric symmetry makes the root
set manageable only when every canonical parent is expanded before insertion.
The resulting candidates are settled by exact deletion intervals and complete
bounded search.  The same deletion theory yields a generalized degree-$(n-1)$
root prepass, while the M\"uhle--Nadeau 3-cycle length gives an exact slack
filter.  Qualifying covers followed by contraction add a one-sided lookup, and
the exhaustive $N=19$ classification adds the degree-18 oriented-unicycle
clearing reduction.  Each reduction has a separate audit or exact integration
test.

Once the complete run and the artifacts in
Table~\ref{tab:final-run} are available, the final paper can replace its
status box and conditional target theorem with the resulting unconditional
diameter statement.

\appendix

\section{Proof-obligation checklist}
\label{app:checklist}

For convenience, the logical dependencies of the negative result are listed
in the order in which a verifier can check them.
\begin{enumerate}
  \item Establish $\TD(17)=10$ and the one-symbol deletion inequality of
    Ara\'ujo et al.~\cite{AraujoEtAl2026}.
  \item Conclude $10\leq\TD(18)\leq11$ and reduce the problem to
    $\E(18,11)$.
  \item Verify the complete $\E(17,10)$ input manifest and its weighted orbit
    count.
  \item Verify the fixed-target extended-toric action, canonical keys, and
    stabilizer weights.
  \item Verify full parent-orbit expansion and all 18 maximum insertions.
  \item Verify deletion standardization, canonical lookup, and complete-store
    absence semantics.
  \item Verify every pruning rule in bounded DFS, including cut-to-$\gamma$
    mapping, the M\"uhle--Nadeau length and cover rules, contractions, inverse
    suppression, slack filtering, both root lookup prepasses, and the $N=19$
    clearing classification.
  \item Verify that all candidates are settled and all parent/generation
    coverage equalities hold.
  \item Verify atomic empty-result publication and absence of cancellation.
  \item Apply Theorem~\ref{thm:certificate}.
\end{enumerate}

\section{Repository map}

The principal implementation and documentation entry points are:
\begin{center}
\begin{tabular}{@{}p{0.48\textwidth}p{0.46\textwidth}@{}}
\toprule
Path & Role\\
\midrule
\path{experiments/README.md} & build, run, storage, GPU, and benchmark guide\\
\path{experiments/docs/TD18_PROOF.md} & concise degree-18 proof checklist\\
\path{experiments/docs/DIAMETRAL_EXTENSION.md} & deletion recurrence\\
\path{experiments/docs/SCAN_BASELINE.md} & independent rank-scan baseline\\
\path{experiments/src/main/java/.../ExtensionRun.java} & extension orchestration\\
\path{experiments/src/main/java/.../ToricOrbit.java} & fixed-target canonicality\\
\path{experiments/src/main/resources/cuda/streaming-count.cu} & exact GPU DFS\\
\bottomrule
\end{tabular}
\end{center}

\bibliographystyle{plain}
\bibliography{references}

\end{document}
